\documentclass[aspectratio=169]{beamer}
\usetheme{metropolis}
\usepackage{graphicx}
\usepackage{tikz}
\usepackage{amsmath}
\usepackage{listings}
\usepackage{xcolor}
\usepackage{booktabs}
\usepackage{emoji}
\usetikzlibrary{shapes,arrows,positioning,calc}

% Code listing setup
\lstset{
  basicstyle=\ttfamily\small,
  keywordstyle=\color{blue},
  commentstyle=\color{gray},
  stringstyle=\color{red},
  showstringspaces=false,
  breaklines=true,
  frame=single,
  language=Python
}

\title{Lecture 12: Attention Mechanisms}
\subtitle{Week 5, Lecture 1 - From Seq2Seq to Attention}
\author{LLM Course}
\date{Winter 2026}

\begin{document}

% ============================================================
% TITLE SLIDE
% ============================================================
\begin{frame}
\titlepage
\end{frame}

% ============================================================
% OVERVIEW
% ============================================================
\begin{frame}{Today's Agenda 📋}
\large
\begin{enumerate}
    \item 🤔 \textbf{The Context Problem}: Why static embeddings aren't enough
    \item 🔄 \textbf{Sequence-to-Sequence Models}: The foundation
    \item ⚠️ \textbf{The Bottleneck Problem}: Why vanilla Seq2Seq struggles
    \item ⚡ \textbf{Attention Mechanisms}: The breakthrough innovation
    \item 🔍 \textbf{How Attention Works}: Step-by-step computation
    \item 👁️ \textbf{Visualizing Attention}: Interpreting the weights
\end{enumerate}

\vspace{1em}
\textit{Goal: Understand why and how attention revolutionized NLP}
\end{frame}

% ============================================================
% THE CONTEXT PROBLEM
% ============================================================
\begin{frame}{The Context Problem 🤔}
\begin{center}
\LARGE
\textbf{Why do we need context-aware models?}
\end{center}

\pause
\vspace{1em}

\begin{exampleblock}{Example: The word "bank"}
\begin{enumerate}
    \item "I deposited money at the \textbf{bank}" (financial institution)
    \item "We sat by the river \textbf{bank}" (riverside)
    \item "The plane started to \textbf{bank} left" (tilt/turn)
\end{enumerate}
\end{exampleblock}

\vspace{1em}

\begin{columns}[T]
\begin{column}{0.48\textwidth}
\textbf{Static Embeddings (Word2Vec):}
\begin{itemize}
    \item One vector per word
    \item \alert{Can't distinguish meanings!}
    \item Context-independent
\end{itemize}
\end{column}

\pause

\begin{column}{0.48\textwidth}
\textbf{Context-Aware Models:}
\begin{itemize}
    \item Different representation per occurrence
    \item Uses surrounding context
    \item Handles polysemy naturally
\end{itemize}
\end{column}
\end{columns}
\end{frame}

% ============================================================
% SEQUENCE TO SEQUENCE INTRO
% ============================================================
\begin{frame}{Sequence-to-Sequence Models 🔄}
\textbf{The breakthrough for variable-length input/output problems}

\vspace{1em}

\textbf{Applications:}
\begin{itemize}
    \item Machine Translation: English → French
    \item Summarization: Long text → Short summary
    \item Question Answering: Question + Context → Answer
    \item Dialogue Systems: User input → System response
\end{itemize}

\vspace{1em}

\begin{center}
\begin{tikzpicture}[scale=0.8]
    % Input sequence
    \foreach \i/\word in {1/The, 2/cat, 3/sat} {
        \node[rectangle,draw,fill=blue!20] at (\i*1.5, 2) {\word};
    }
    \node at (0, 2) {Input:};

    % Arrow
    \draw[->, very thick] (5, 2) -- (6.5, 2);
    \node[rectangle,draw,fill=green!20] at (7.5, 2) {Model};
    \draw[->, very thick] (8.5, 2) -- (10, 2);

    % Output sequence
    \node at (10.5, 2) {Output:};
    \foreach \i/\word in {1/Le, 2/chat, 3/s'assit} {
        \node[rectangle,draw,fill=red!20] at (10.5+\i*1.5, 2) {\word};
    }
\end{tikzpicture}
\end{center}

\vspace{1em}
\small
\textit{Reference: Sutskever et al. (2014) - "Sequence to Sequence Learning with Neural Networks"}
\end{frame}

% ============================================================
% ENCODER DECODER ARCHITECTURE
% ============================================================
\begin{frame}{Encoder-Decoder Architecture 🏗️}
\textbf{Two-part architecture: Encode then Decode}

\vspace{1em}

\begin{center}
\begin{tikzpicture}[
    node distance=0.8cm,
    rnn/.style={rectangle,draw,fill=blue!20,minimum width=1cm,minimum height=1cm},
    decoder/.style={rectangle,draw,fill=red!20,minimum width=1cm,minimum height=1cm}
]
    % Encoder
    \node[rnn] (e1) {RNN};
    \node[rnn,right=of e1] (e2) {RNN};
    \node[rnn,right=of e2] (e3) {RNN};

    % Input words
    \node[below=of e1] {The};
    \node[below=of e2] {cat};
    \node[below=of e3] {sat};

    % Context vector
    \node[rectangle,draw,fill=green!20,minimum width=1.5cm,minimum height=1cm,right=of e3] (ctx) {Context\\Vector};

    % Decoder
    \node[decoder,right=of ctx] (d1) {RNN};
    \node[decoder,right=of d1] (d2) {RNN};
    \node[decoder,right=of d2] (d3) {RNN};

    % Output words
    \node[above=of d1] {Le};
    \node[above=of d2] {chat};
    \node[above=of d3] {s'assit};

    % Arrows
    \draw[->] (e1) -- (e2);
    \draw[->] (e2) -- (e3);
    \draw[->] (e3) -- (ctx);
    \draw[->] (ctx) -- (d1);
    \draw[->] (d1) -- (d2);
    \draw[->] (d2) -- (d3);

    % Labels
    \node[above=2cm of e2,blue] {\textbf{Encoder}};
    \node[above=2cm of d2,red] {\textbf{Decoder}};
\end{tikzpicture}
\end{center}

\vspace{1em}

\begin{columns}[T]
\begin{column}{0.48\textwidth}
\textbf{Encoder:}
\begin{itemize}
    \item Reads input sequence
    \item Compresses to fixed-size vector
    \item Captures semantic meaning
\end{itemize}
\end{column}

\begin{column}{0.48\textwidth}
\textbf{Decoder:}
\begin{itemize}
    \item Starts from context vector
    \item Generates output sequence
    \item One token at a time
\end{itemize}
\end{column}
\end{columns}
\end{frame}

% ============================================================
% SEQ2SEQ PROBLEM
% ============================================================
\begin{frame}{The Seq2Seq Bottleneck Problem ⚠️}
\textbf{Challenge: All information compressed into single vector!}

\vspace{1em}

\begin{alertblock}{The Problem}
\begin{itemize}
    \item Long sequences → information loss
    \item Fixed-size context vector is a bottleneck
    \item Early tokens forgotten by the time we reach the end
    \item Performance degrades with sequence length
\end{itemize}
\end{alertblock}

\vspace{1em}

\begin{center}
\begin{tikzpicture}[scale=0.9]
    % Long input
    \foreach \i in {1,...,8} {
        \node[rectangle,draw,fill=blue!20,minimum size=0.5cm] at (\i*0.7, 2) {};
    }
    \node at (0, 2) {Input:};

    % Bottleneck
    \draw[->] (6.5, 2) -- (7.5, 2);
    \node[rectangle,draw,fill=orange!50,minimum width=0.8cm,minimum height=1.5cm] at (8.5, 2) {};
    \node[below=0.1cm of {(8.5,2)}] {\alert{Bottleneck!}};
    \draw[->] (9.5, 2) -- (10.5, 2);

    % Output
    \foreach \i in {1,...,8} {
        \node[rectangle,draw,fill=red!20,minimum size=0.5cm] at (10.5+\i*0.7, 2) {};
    }
\end{tikzpicture}
\end{center}

\vspace{1em}

\begin{center}
\LARGE
\textbf{Solution: Attention! ⚡}
\end{center}
\end{frame}

% ============================================================
% ATTENTION MECHANISM
% ============================================================
\begin{frame}{Attention Mechanism: The Big Idea ⚡}
\textbf{Instead of compressing everything into one vector...}

\vspace{0.5em}

\textbf{Let the decoder look at all encoder hidden states!}

\vspace{1em}

\begin{center}
\begin{tikzpicture}[
    node distance=0.5cm,
    enc/.style={rectangle,draw,fill=blue!20,minimum width=0.8cm,minimum height=0.8cm},
    dec/.style={rectangle,draw,fill=red!20,minimum width=0.8cm,minimum height=0.8cm}
]
    % Encoder states
    \node[enc] (e1) at (0, 0) {$h_1$};
    \node[enc] (e2) at (1.2, 0) {$h_2$};
    \node[enc] (e3) at (2.4, 0) {$h_3$};
    \node[enc] (e4) at (3.6, 0) {$h_4$};

    \node[left=0.3cm of e1] {Encoder:};

    % Decoder state
    \node[dec] (d1) at (1.8, 3) {$s_t$};
    \node[left=0.3cm of d1] {Decoder:};

    % Attention arrows
    \draw[->, thick, blue] (e1) -- (d1) node[midway,left] {\tiny 0.1};
    \draw[->, thick, orange, line width=2pt] (e2) -- (d1) node[midway,left] {\tiny \textbf{0.6}};
    \draw[->, thick, orange, line width=2pt] (e3) -- (d1) node[midway,right] {\tiny \textbf{0.25}};
    \draw[->, thick, blue] (e4) -- (d1) node[midway,right] {\tiny 0.05};

    \node[below=1.5cm of e2] {Input: "The cat sat down"};
    \node[above=0.5cm of d1] {Generating: "Le \textbf{chat}"};
    \node[right=2cm of d1] {\alert{Attention weights sum to 1.0}};
\end{tikzpicture}
\end{center}

\vspace{1em}

\textbf{Key Insight:} When generating "chat" (cat), pay more attention to "cat" in the input!

\vspace{0.5em}
\small
\textit{Reference: Bahdanau et al. (2015) - "Neural Machine Translation by Jointly Learning to Align and Translate"}
\end{frame}

% ============================================================
% ATTENTION COMPUTATION
% ============================================================
\begin{frame}{How Attention Works: Step by Step 🔍}
\textbf{Computing attention weights:}

\vspace{1em}

\begin{enumerate}
    \item \textbf{Score}: How relevant is each encoder state to current decoder state?
    \begin{align*}
    e_{t,i} = \text{score}(s_t, h_i) = s_t^T W_a h_i
    \end{align*}

    \item \textbf{Normalize}: Convert scores to probabilities (softmax)
    \begin{align*}
    \alpha_{t,i} = \frac{\exp(e_{t,i})}{\sum_{j=1}^{T} \exp(e_{t,j})}
    \end{align*}

    \item \textbf{Context}: Weighted sum of encoder states
    \begin{align*}
    c_t = \sum_{i=1}^{T} \alpha_{t,i} h_i
    \end{align*}

    \item \textbf{Decode}: Use context vector along with decoder state
    \begin{align*}
    s_{t+1} = f(s_t, c_t, y_t)
    \end{align*}
\end{enumerate}

\vspace{0.5em}

\textbf{Result:} Decoder dynamically focuses on different parts of input!
\end{frame}

% ============================================================
% ATTENTION TYPES
% ============================================================
\begin{frame}{Attention Score Functions 📐}
\textbf{Different ways to compute the score}

\vspace{1em}

\begin{enumerate}
    \item \textbf{Additive Attention (Bahdanau et al. 2015):}
    \begin{align*}
    \text{score}(s_t, h_i) = v^T \tanh(W_1 s_t + W_2 h_i)
    \end{align*}
    \begin{itemize}
        \item Learns alignment jointly with translation
        \item More parameters, more flexibility
    \end{itemize}

    \vspace{0.5em}

    \item \textbf{Multiplicative Attention (Luong et al. 2015):}
    \begin{align*}
    \text{score}(s_t, h_i) = s_t^T W h_i
    \end{align*}
    \begin{itemize}
        \item Simpler, faster computation
        \item Works well in practice
    \end{itemize}

    \vspace{0.5em}

    \item \textbf{Scaled Dot-Product (Vaswani et al. 2017):}
    \begin{align*}
    \text{score}(s_t, h_i) = \frac{s_t^T h_i}{\sqrt{d}}
    \end{align*}
    \begin{itemize}
        \item Used in Transformers (coming next lecture!)
        \item Scaling prevents vanishing gradients
    \end{itemize}
\end{enumerate}

\small
\textit{Reference: Luong et al. (2015) - "Effective Approaches to Attention-based Neural Machine Translation"}
\end{frame}

% ============================================================
% ATTENTION VISUALIZATION
% ============================================================
\begin{frame}{Attention Visualization 👁️}
\textbf{Example: English → French translation with attention weights}

\vspace{1em}

\begin{center}
\begin{tabular}{c|cccc}
\textbf{Output} & \multicolumn{4}{c}{\textbf{Input Words}} \\
\textbf{(French)} & The & agreement & on & European \\
\hline
L' & \cellcolor{orange!80}0.7 & 0.1 & 0.1 & 0.1 \\
accord & 0.1 & \cellcolor{orange!80}0.8 & 0.05 & 0.05 \\
sur & 0.05 & 0.05 & \cellcolor{orange!80}0.8 & 0.1 \\
l'européen & 0.05 & 0.05 & 0.1 & \cellcolor{orange!80}0.8 \\
\end{tabular}
\end{center}

\vspace{1em}

\textbf{Observations:}
\begin{itemize}
    \item Diagonal pattern for similar word order
    \item Model learns alignment automatically!
    \item No need for explicit word alignment annotations
    \item Can handle reordering (syntax differences between languages)
\end{itemize}

\vspace{0.5em}
\textit{Attention provides interpretability: we can see what the model is "looking at"}
\end{frame}

% ============================================================
% ATTENTION BENEFITS
% ============================================================
\begin{frame}{Benefits of Attention Mechanisms ✅}
\begin{enumerate}
    \item \textbf{Solves the Bottleneck Problem}
    \begin{itemize}
        \item Decoder has access to all encoder states
        \item No information compression into single vector
        \item Works well for long sequences
    \end{itemize}

    \vspace{0.5em}

    \item \textbf{Improves Performance}
    \begin{itemize}
        \item Better BLEU scores on translation tasks
        \item Handles long-range dependencies
        \item More robust to sequence length
    \end{itemize}

    \vspace{0.5em}

    \item \textbf{Provides Interpretability}
    \begin{itemize}
        \item Can visualize what the model focuses on
        \item Helps debug and understand model behavior
        \item Builds trust in model predictions
    \end{itemize}

    \vspace{0.5em}

    \item \textbf{Enables Better Alignment}
    \begin{itemize}
        \item Learns source-target correspondences
        \item No need for external alignment tools
        \item Works across different language pairs
    \end{itemize}
\end{enumerate}

\vspace{0.5em}
\textbf{Attention became the foundation for modern NLP!}
\end{frame}

% ============================================================
% REAL WORLD IMPACT
% ============================================================
\begin{frame}{Real-World Impact of Attention 🌍}
\textbf{Attention mechanisms revolutionized multiple domains:}

\vspace{1em}

\begin{columns}[T]
\begin{column}{0.48\textwidth}
\textbf{Machine Translation:}
\begin{itemize}
    \item Google Translate (2016)
    \item DeepL
    \item Facebook translations
    \item Dramatic quality improvements
\end{itemize}

\vspace{0.5em}

\textbf{Text Summarization:}
\begin{itemize}
    \item News article summarization
    \item Document understanding
    \item Email auto-responses
\end{itemize}
\end{column}

\begin{column}{0.48\textwidth}
\textbf{Question Answering:}
\begin{itemize}
    \item Reading comprehension
    \item Search engines
    \item Virtual assistants
\end{itemize}

\vspace{0.5em}

\textbf{Speech Recognition:}
\begin{itemize}
    \item Attend to acoustic features
    \item Better transcription accuracy
    \item Listen, Attend and Spell models
\end{itemize}
\end{column}
\end{columns}

\vspace{1em}

\begin{block}{Key Milestone}
By 2016, attention-based models became the standard for sequence-to-sequence tasks, paving the way for the Transformer revolution in 2017.
\end{block}
\end{frame}

% ============================================================
% CODE EXAMPLE
% ============================================================
\begin{frame}[fragile]{Implementing Attention in PyTorch 💻}
\textbf{Simple attention mechanism implementation}

\begin{lstlisting}[language=Python]
import torch
import torch.nn as nn
import torch.nn.functional as F

class BahdanauAttention(nn.Module):
    def __init__(self, hidden_dim):
        super().__init__()
        self.W_dec = nn.Linear(hidden_dim, hidden_dim)
        self.W_enc = nn.Linear(hidden_dim, hidden_dim)
        self.v = nn.Linear(hidden_dim, 1)

    def forward(self, decoder_hidden, encoder_outputs):
        # decoder_hidden: [batch, hidden_dim]
        # encoder_outputs: [batch, seq_len, hidden_dim]

        # Compute scores
        dec = self.W_dec(decoder_hidden).unsqueeze(1)  # [batch, 1, hidden]
        enc = self.W_enc(encoder_outputs)  # [batch, seq_len, hidden]
        scores = self.v(torch.tanh(dec + enc))  # [batch, seq_len, 1]

        # Compute attention weights
        attn_weights = F.softmax(scores, dim=1)  # [batch, seq_len, 1]

        # Compute context vector
        context = torch.sum(attn_weights * encoder_outputs, dim=1)

        return context, attn_weights.squeeze(-1)
\end{lstlisting}
\end{frame}

% ============================================================
% DISCUSSION
% ============================================================
\begin{frame}{Discussion Questions 💭}
\begin{enumerate}
    \item \textbf{Why is attention called "soft alignment"?}
    \begin{itemize}
        \item How is it different from hard alignment?
        \item What are the advantages of soft vs. hard?
    \end{itemize}

    \vspace{0.5em}

    \item \textbf{Computational Cost:}
    \begin{itemize}
        \item What is the time complexity of attention?
        \item How does it scale with sequence length?
        \item When might this be a problem?
    \end{itemize}

    \vspace{0.5em}

    \item \textbf{Interpretability:}
    \begin{itemize}
        \item Can we always trust attention weights as explanations?
        \item What about when attention is uniform across all inputs?
    \end{itemize}

    \vspace{0.5em}

    \item \textbf{Beyond Seq2Seq:}
    \begin{itemize}
        \item Where else might attention be useful?
        \item Can we apply it within a single sequence?
        \item (Spoiler: Yes! That's self-attention → next lecture!)
    \end{itemize}
\end{enumerate}
\end{frame}

% ============================================================
% LOOKING AHEAD
% ============================================================
\begin{frame}{Looking Ahead 🔮}
\textbf{What's Next?}

\vspace{1em}

\textbf{Today we learned:}
\begin{itemize}
    \item The context problem in NLP
    \item Sequence-to-sequence architecture
    \item The bottleneck problem
    \item How attention mechanisms work
    \item Attention as alignment and interpretation
\end{itemize}

\vspace{1em}

\textbf{Next lecture (Lecture 13):}
\begin{itemize}
    \item \alert{Self-Attention}: Attention within a sequence
    \item \alert{The Transformer}: "Attention is All You Need"
    \item \alert{Query, Key, Value}: The new framework
    \item \alert{Multi-Head Attention}: Learning diverse relationships
    \item \alert{Positional Encoding}: Injecting order information
\end{itemize}

\vspace{1em}

\begin{center}
\Large
\textbf{Get ready for the Transformer revolution! 🚀}
\end{center}
\end{frame}

% ============================================================
% SUMMARY
% ============================================================
\begin{frame}{Summary 🎯}
\textbf{Key Takeaways:}

\vspace{1em}

\begin{enumerate}
    \item \textbf{Context Matters}
    \begin{itemize}
        \item Static embeddings can't capture context-dependent meanings
        \item Need dynamic representations based on context
    \end{itemize}

    \item \textbf{Seq2Seq Bottleneck}
    \begin{itemize}
        \item Fixed-size context vector limits performance
        \item Information loss for long sequences
    \end{itemize}

    \item \textbf{Attention is the Solution}
    \begin{itemize}
        \item Dynamic access to all encoder states
        \item Weighted combination based on relevance
        \item Attention weights sum to 1.0 (probability distribution)
    \end{itemize}

    \item \textbf{Benefits}
    \begin{itemize}
        \item Better performance on long sequences
        \item Automatic alignment learning
        \item Model interpretability
    \end{itemize}
\end{enumerate}

\vspace{0.5em}
\textbf{Attention mechanisms laid the foundation for modern transformer-based models!}
\end{frame}

% ============================================================
% REFERENCES
% ============================================================
\begin{frame}{References 📚}
\textbf{Key Papers:}

\vspace{0.5em}

\begin{itemize}
    \item \textbf{Sutskever et al. (2014)} - "Sequence to Sequence Learning with Neural Networks"
    \begin{itemize}
        \item Introduced encoder-decoder architecture
        \item Foundation for seq2seq models
    \end{itemize}

    \vspace{0.3em}

    \item \textbf{Bahdanau et al. (2015)} - "Neural Machine Translation by Jointly Learning to Align and Translate"
    \begin{itemize}
        \item Introduced additive attention mechanism
        \item Solved the bottleneck problem
    \end{itemize}

    \vspace{0.3em}

    \item \textbf{Luong et al. (2015)} - "Effective Approaches to Attention-based Neural Machine Translation"
    \begin{itemize}
        \item Multiplicative attention variants
        \item Global vs. local attention
    \end{itemize}

    \vspace{0.3em}

    \item \textbf{Vaswani et al. (2017)} - "Attention Is All You Need"
    \begin{itemize}
        \item The Transformer architecture (next lecture!)
        \item Scaled dot-product attention
    \end{itemize}
\end{itemize}

\vspace{0.5em}

\textbf{Additional Resources:}
\begin{itemize}
    \item Jay Alammar's blog: "Visualizing A Neural Machine Translation Model"
    \item Distill.pub: "Attention and Augmented Recurrent Neural Networks"
\end{itemize}
\end{frame}

% ============================================================
% QUESTIONS
% ============================================================
\begin{frame}{Questions? 🙋}
\begin{center}
\LARGE
\textbf{Discussion Time}
\end{center}

\vspace{2em}

\textbf{Office Hours Topics:}
\begin{itemize}
    \item Implementing attention from scratch
    \item Different attention mechanisms
    \item Debugging attention-based models
    \item Assignment 4 preparation
\end{itemize}

\vspace{2em}

\begin{center}
\Large
Thank you! 🙏

\vspace{1em}

See you next lecture for Transformers!
\end{center}
\end{frame}

\end{document}
